\section{局部几何朗兰兹与仿射 Beilinson--Bernstein 局部化}

\label{s:affbbloc}

\subsection{概览}

\ssszh 前几节讨论了约化李代数的局部化理论，并说明了它与约化群的范畴表示论之间的关系。第 \ref{s:catloop} 节中，我们遇到了环路群及其范畴表示。

在这最后一节里，我们想概述局部几何朗兰兹中的一些基本猜想与结果；它们可以理解为对环路群的范畴表示作参数化。最后我们会回到原点，遇到环路群上 Beilinson--Bernstein 局部化的类比，它由 Beilinson--Drinfeld、Frenkel--Gaitsgory 与 Raskin--Yang 发展起来。

\subsection{对应的朴素表述}



\ssszh 对 $p$-adic 群 $G(\Fqt)$，局部朗兰兹对应（猜想地）用 Galois 理论的数据来参数化它在 $k$ 向量空间上的不可约表示；初步地说，这些数据是从 $\Fqt$ 的绝对 Galois 群到朗兰兹对偶 $\check{G}(k)$ 的同态。


本小节想综述：对 $\Dmod(LG)\mod$ 的类似参数化，人们期待什么、又已知什么。



\ssszh 回忆 $k(\!(t)\!)$ 由形式穿孔圆盘 $D^\times$ 上的函数构成。按照 Galois 表示与局部系统之间的类比（参见附录 \ref{s:frmfun}），当前设定中与 $\cG$-值 Galois 表示相对应的，是平坦的 $\check{G}$-联络，即 $D^\times$ 上的 $\cG$ de Rham 局部系统。

具体地说，它们由联络模去规范等价给出，即
%
\[
    \LS \simeq \{ d + \cfg(\!(t)\!)dt \} / \cG(k(\!(t)\!)).
\]
%
同一个公式还把 $\LS$ 定义为代数几何的对象：取 $L\cG$ 在联络形式 $d + \cfg(\!(t)\!)dt$ 的归纳概形上作用的商预叠，而后者与例 \ref{ex:ga} 类似地定义；进一步的讨论参见 \cite{raskinspectraldecomp}。



\ssszh \label{s:infdiscuss} 下面试着以略为启发式的方式给出局部几何朗兰兹对应的陈述。初步地说，可以期望不可约的 $\Dmod(LG)$ 表示与 $\LS$ 的点一一对应。也就是说，对每个 $\cG$-联络 $\sigma$，应当能典范地配上 $LG$ 的一个不可约范畴表示 $\sC_{\sigma}$。

由于 $\Dmod(LG)\mod$ 的一般对象 $\sC$ 应当分解为不可约表示的直积分，它应当定义 $\LS$ 上相应的一族范畴，而该族在每个点的纤维衡量 $\sC_{\sigma}$ 在直积分中相应的重数。

由于不可约性是阿贝尔范畴中对象的特征，这段讨论很难精确表述。不过，对于空间上的一族范畴，有一个好的概念，即拟凝聚的范畴层，下面就来综述。

作一些修正（下文公式 \eqref{e:loclang} 会适时遇到）之后，这就给出局部几何朗兰兹对应的陈述。也就是说，猜想是：环路群的全部范畴表示与局部系统模空间上的范畴层之间，应当有 $(\infty,2)$-范畴的等价，即
%
\begin{equation}
\label{e:loclangnaive}       \mathbb{L}: \Dmod(LG)\mod \simeq
\on{2-QCoh}(\LS).
\end{equation}
%



下面给出一条帮助定位的基本注。所述 2-范畴的等价并不是说：对给定的范畴表示 $\sC$，$\sC$ 与 $\mathbb{L}(\sC)$ 的底层 dg 范畴等价。而是说：对任意一对对象 $\sC$ 与 $\sD$，有等价
%
\[
         \Hom_{\Dmod(LG)\mod}(\sC, \sD) \simeq
         \Hom_{\TwoCoh(\LS)}(\mathbb{L}(\sC), \mathbb{L}(\sD)),
\]
%
并且当 $\sC$、$\sD$ 变动时，这些等价与复合相容。

\begin{re}  作为玩具模型，通常代数的 Morita 等价
    \[
        A\mod \simeq B\mod
    \]
	%
	匹配的并不是模的底层向量空间，而只是它们之间的同态。具体地说，对有限维 $k$ 向量空间 $V$，可以考察 Morita 等价
	%
	\[
	   V \otimes -  :\Vect \simeq  \End(V)\mod.
	\]
	%
	它把 $k$ 与 $V$ 互换，而这两者当然不同构；但在自同态上确实有
	%
	\[
	     \Hom_{\Vect}(k, k) \simeq k \simeq \Hom_{\End(V)}(V,V).
	\]
	%
		\end{re}


在转向修正之前，我们想先讨论范畴层，以及 $\mathbb{L}$ 应当满足的一些基本相容性。

\ssszh   要讨论范畴层，先从仿射簇 $X = \on{Spec }O_X$ 的情形开始。此时，$X$ 上的拟凝聚范畴层不过是一个 $O_X$-线性的 dg 范畴（而非仅仅是 $k$-线性的）。由于 $O_X$-线性等同于可以张量 $O_X$-模的复形，等价地有
%
\[
	  \TwoCoh(X) \simeq \QCoh(X)\mod.
\]
%
特别地，$\TwoCoh(\on{pt})$ 就是 $\DGCat$；并且对任意闭点 $\on{pt} \rightarrow X$，有拉回
%
\[
	  \Vect \otimes_{\QCoh(X)}  -: \TwoCoh(X) \rightarrow \DGCat,
\]
%
也就是说，可以谈论范畴层的纤维，它们就是通常的 dg 范畴。这样，范畴层就被看成 $X$ 上的一族范畴。



\begin{re}\label{r:pushpull2coh} 更一般地，给定任意映射 $f: X \rightarrow Y$，拉回定义了幺半函子 $\QCoh(Y) \rightarrow \QCoh(X)$，特别地，得到伴随的诱导函子与限制函子
	%
	\[
	    f^*: \TwoCoh(Y) \leftrightarrows \TwoCoh(X): f_*.
	\]
	%
\end{re}





下面是一个范畴层的基本例子。

\begin{ex}给定簇的映射 $\tilde{X} \rightarrow X$，$\QCoh(\tilde{X})$ 是 $\QCoh(X)$-代数，特别地是 $X$ 上的范畴层。对 $X$ 的任意点 $x$，由例 \ref{e:fibprodvar}，该范畴层的纤维不过是实际纤维上的拟凝聚层，即
%
\[
	  i_x^* \QCoh(\tilde{X}) \simeq \QCoh( \tilde{X} \times_X x).
\]
%
\end{ex}


\ssszh 对一般的 $X$（不必仿射），定义像拟凝聚层的定义那样通过粘接给出。细节对我们并不那么重要，但应当提到 Gaitsgory 在 \cite{1affine} 中建立的下述漂亮现象。即：对一般的 $X$，按定义是逐个仿射地拼贴范畴层；另一方面，又有拟凝聚层的整体幺半范畴，从而有其左模范畴 $\QCoh(X)\mod$。于是可以问这两者之间的关系。与通常的拟凝聚层一样，有整体截面函子
%
\begin{equation} \label{e:globsectshvcat}
 \Gamma: \TwoCoh(X) \rightarrow \QCoh(X)\mod.
\end{equation}
%
对通常的层，类似的映射
    %
    \[
    \Gamma: \QCoh(X) \rightarrow \Gamma(X, \OO_X)\mod
    \]
	%
	在 $X$ 非仿射时极少是等价。\footnote{对「极少」一词有疑问的读者：要点在于，若考虑整体函数的完整导出代数，则对非仿射簇（例如仿射簇的开子簇），$\Gamma$ 仍可以是同构。}
相反，对多数有限维的 $X$（例如任意有限型的代数叠），函子 \eqref{e:globsectshvcat} 都是等价；这一现象称作 1-仿射性。


本领域的一条基本猜想（属于 Gaitsgory）是：$\LS$ 也是 1-仿射的，参见 \cite{raskinspectraldecomp}。这是一条微妙的断言，因为对归纳概形或无限型群概形的分类空间，1-仿射性一般并不成立。所以从某种意义上说，需要这两类问题对 $\LS$ 恰好「抵消」。Raskin 的一项重要工作 \cite{raskinspectraldecomp} 部分地处理了这一点，他证明了：对 $X = \LS$，\eqref{e:globsectshvcat} 是全忠实的。
